\section{A growing mixed cubic family on the original window}
\label{sec:mixed-cubic-split-prime-window}

\paragraph{The preceding original-coefficient remainder.}
Retain every datum, page, prescribed pair, field, original coefficient,
lattice, normalization and full local condition of AC-G202-window.
In particular $h\geq2$ and $T\geq2$ are integers, and at source size $u$
\[
 \begin{gathered}
 H_u=\mathbb Q(\zeta_{2^{h+2}},\zeta_7,\sqrt{-q},\sqrt{-d},
                  \omega,\sqrt[3]{269}),\qquad K=\mathbb Q(\omega),\\
 u^{2+8\varepsilon}\leq2^h\leq2u^{2+8\varepsilon},\qquad
 T=\lfloor u^{399}\rfloor,\qquad0<\varepsilon<1/256.
 \end{gathered}\tag{MCW.1}
\]
Let $L=K(\sqrt[3]{269})$ and let $\chi$ be either faithful
character of $\operatorname{Gal}(L/K)=C_3$.
Keep exactly the complete preceding omitted set $Z_u$, its
retained partition $J_u=J_{u,0}\sqcup J_{u,\rm mix}$, and the
sets $V_u,Y_u,C_5(u),U_5(u),C_7^{\rm all}(u),U_7^{\rm all}(u)$
of (CAW.1)--(CAW.4), (F5W.2) and (S7A.2).
The preceding explicit remainder is (S7A.3).

Define the same explicit prime criterion as in (MCP.1):
\[
 \begin{split}
 \mathscr P_{\rm cub}=\{\ell\text{ prime}:{}&\ell>3,\quad
       \ell\ne7,269,\quad\ell\equiv1\pmod3,\\
       &269^{(\ell-1)/3}\not\equiv1\pmod\ell\},\\
 M(u)=\{(\ell,\chi^j,e):{}&\ell\in\mathscr P_{\rm cub},\quad
                 j=1,2,\quad1\leq e<T,\quad(\ell-1)\mid e\},\\
 J_{\rm cub}(T)&=\sum_{\substack{\ell\in\mathscr P_{\rm cub}\\\ell\leq T}}
                     \left\lfloor\frac{T-1}{\ell-1}\right\rfloor.
 \end{split}\tag{MCW.2}
\]
Every undisplayed rational two-power, conductor-seven and auxiliary
$q,d$ factor of a label in $M(u)$ is trivial. No finite
Teichm\"uller character at $\ell$ is added to the host.

Define
\[
 \begin{split}
 \mathcal R_{399}^{\rm match,cub}(u)={}&
 \sum_{(\ell,\eta,e)\in J_{u,\rm mix}\setminus M(u)}
                     \frac{\log\#\mathcal S_{\ell,\eta}(e)}{b_\ell}\\
 &+\sum_{(\ell,\eta,e)\in
     J_{u,0}\setminus(V_u\cup Y_u\cup C_7^{\rm all}(u)\cup C_5(u))}
       \frac{\log\#\mathcal S^{\mathbb Q}_{\ell,
                                \psi_{\eta,e}^{\rm nc}}(e)}{b_\ell}.
 \end{split}\tag{MCW.3}
\]
The first sum uses the full original $K$-strict groups. The second
rational noncritical sum is unchanged from (S7A.3).
Finiteness is the one already proved in AC-G191/AC-G192.

\begin{theorem}[AC-G204-window: mixed cubic deletion at varying coefficient primes]
\label{AC-G204-window}
Every triple of $M(u)$ is an actual original-coefficient strict
zero, and the following inclusions and equalities hold:
\[
 \begin{gathered}
 M(u)\subset J_{u,\rm mix},\qquad
 M(u)\cap\bigl(Z_u\cup V_u\cup Y_u\cup C_7^{\rm all}(u)\cup C_5(u)\bigr)
                      =\varnothing,\\
 |M(u)|=2J_{\rm cub}(T),\qquad
 \mathcal R_{399}^{\rm match,cub}(u)
                        =\mathcal R_{399}^{\rm match,7all,5}(u).
 \end{gathered}\tag{MCW.4}
\]
This is exactly $2J_{\rm cub}(T)$ newly absent mixed terms,
without subtracting an old overlap. In particular
\[
 |M(u)|\geq
  2\left\lfloor\frac{T-1}{12}\right\rfloor+
  2\left\lfloor\frac{T-1}{18}\right\rfloor.
 \tag{MCW.5}
\]
For $T\geq13$ the retained indexing set is properly reduced.
The original triples no longer represented form the disjoint union
\[
 Z_u\sqcup V_u\sqcup Y_u\sqcup U_7^{\rm all}(u)\sqcup U_5(u)\sqcup M(u).
 \tag{MCW.6}
\]
Its cardinality is the exact (S7A.6) count plus $2J_{\rm cub}(T)$.
No individual positive derived index is asserted zero.

With the same original constants and costs as (S7A.7),
\[
 \begin{aligned}
 |I_{399}-\mathcal R_{399}^{\rm match,cub}|&\leq
        \mathcal E_u+\mathcal L_u+\mathcal A_u+C_2\mathcal P_u,\\
 |I_{399}^{\rm all7,2}-\mathcal R_{399}^{\rm match,cub}|&\leq
        \mathcal E_u+\mathcal L_u+C_2\mathcal P_u.
 \end{aligned}\tag{MCW.7}
\]
No additional local, torsion, critical or prime-count cost appears.
The envelope remains $O_\varepsilon(u^{806+16\varepsilon}(\log u)^2)$.
The preceding and new numerical remainder premises use the
identical window, threshold and slack. The original full index
premise follows by the same threshold increase and half-slack
argument, with all original simultaneous quantifiers unchanged.
\end{theorem}

\begin{proof}
\proofstep{1. Actual mixed characters and the literal original-ring theorem.}
The exact product decomposition (FC.1) has the direct factor
$\operatorname{Gal}(L/K)=C_3$. Its two nontrivial characters,
with every other factor trivial, are the literal labels in (MCW.2).
Their values belong to the prescribed original completion of
$\mathbb Z[\zeta_{3\cdot2^h}]$. In fact $\ell\equiv1\pmod3$
already puts these roots in $\mathbb Q_\ell$, but the original
larger $O_\ell$ and its chosen embedding are retained.
The coefficient is canonically
\[
 (O_\ell\otimes_{\mathbb Z_\ell}\mathbb Q_\ell/\mathbb Z_\ell)
            (\chi^j)(-e)
       =(E_\ell/O_\ell)(\chi^j c_\ell^{-e}).
 \tag{MCW.8}
\]
Apply AC-G203 (MCP.3) and (MCP.18) with $\kappa=1$.
The hypotheses on $\ell$ and $e$ are precisely (MCW.2).
This gives zero for the full original group in the first sum
of (MCW.3), including both completions at the split coefficient
prime and every away-inertia condition.
The argument in G203 verifies residual Frobenius nontriviality
at both places, the exact ordinary class-quotient input, and the
original-uniformizer lift. No new rank-one local generation or
unit-basis premise is substituted for it.

This removes a full mixed $K$-strict group. It does not pass
through a rational parity component or add a critical partner.
The original $b_\ell$, ramification, codifferent and finite
localization are unchanged. The finite locally constant etale
coefficients on $\operatorname{Spec}\mathcal O_K[1/(3\cdot269\cdot\ell)]$
and full completion/inertia dictionary are exactly the retained
finite-cover interpretation in G203 and (FC.2).

\proofstep{2. Exhaustive old-overlap audit.}
Every triple of $M(u)$ has nontrivial Kummer cubic factor.
All earlier fixed-coefficient mixed omissions in $Z_u$ have
coefficient $269$, excluded in (MCW.2).
Outside that coefficient, the other old fixed families have
coefficients two or three, also excluded. The only remaining
old $Z_u$ possibilities are the high-two families (OPW.3) and
the inert-trivial families (I2W.3); both have trivial Kummer
cubic factor. Thus no member of $M(u)$ lies in $Z_u$.
This conclusion does not identify the cubic part with a
rational modulo-seven character of order three: they are
independent factors in the exact original group.

The sets $V_u,Y_u,C_5(u),C_7^{\rm all}(u)$ likewise have trivial
Kummer cubic part. This stays true if a new coefficient prime
happens to equal an auxiliary prime $q$ or $d$: its old $Y_q$
or $Y_d$ row has a rational Legendre factor and trivial Kummer
factor, whereas the new label has the opposite factor choices.
The direct product makes these distinct.
This proves the first line of (MCW.4) and the disjoint union
(MCW.6). Every other mixed or rational noncritical row remains.

\proofstep{3. Exact growing-prime count and pointwise deletion.}
For each eligible $\ell$, the positive weights are
$(\ell-1),2(\ell-1),\ldots$ below $T$, giving the summand
$\lfloor(T-1)/(\ell-1)\rfloor$. If it is nonzero then
$\ell\leq T$, so (MCW.2) is a finite sum. This upper bound
arises from the weight condition; it is not a replacement of
the unrestricted coefficient-prime domain in (TC.1).
The two faithful cubic characters are distinct on the direct
$C_3$ factor. Primes and weights are separate coordinates.
This proves the exact count in (MCW.4).

The explicit checks $269\equiv9\pmod{13}$, $9^4\equiv9\ne1$,
and $269\equiv3\pmod{19}$, $3^6\equiv7\ne1$, show that both
primes are eligible. Their weight progressions yield (MCW.5).
At $T\geq13$, $(13,\chi,12)$ is a newly absent triple.
No density theorem or unproved quantitative prime estimate is
needed for this properness or the exact count.

Delete the zero mixed summands supplied in Step1, preserving
every normalization. Equivalently apply (MCP.6) to the retained
set after all preceding omissions. This proves the final equality
of (MCW.4) pointwise at the same source size $u$.
The rational noncritical part and its critical comparison do
not change. No positive derived determinant is split or set to zero.

\proofstep{4. The inherited single cost and common window.}
Substitute the equality (MCW.4) into (S7A.7).
This proves (MCW.7) with one full (FL.4) local majorant,
one (LT.2) torsion aggregate, one existing critical aggregate
and the same G189 cost against the full index.
There is no multiplier for $|\mathscr P_{\rm cub}\cap[1,T]|$,
$J_{\rm cub}(T)$ or the number of weights.
The ambient degree $D_u=144\cdot2^h$ and abelian degree
$d_u=96\cdot2^h$ retain their distinct source roles.

Fix the original datum, page and $0<\varepsilon<1/256$.
Retain the exact (TC.3) parameters $J,X_0,s$, prescribed pairs
and the same $X^{1/J}\leq W$, $W^2\leq X$.
Equal explicit remainders have the identical numerical premise
on that same window, with the same threshold and slack.
For the original full index premise take the same uniform $U_0$
that makes the first error in (MCW.7) at most
$(s/2)\exp(u^\varepsilon/2)$, increase the threshold only to
$\max(X_0,U_0^J)$, and retain that same $W$ with slack $s/2$.
No independent existential windows are intersected.
The actual source rank $R.t$ and every simultaneous field,
coefficient, lattice, page, pair and finite condition remain.
\end{proof}

\paragraph{Scope and chronology.}
The prime criterion and full mixed strict-zero theorem are the
completed AC-G203 input; this composition supplies its exhaustive
old-overlap and identical-window application.
The exact arbitrary-set deletion identity was already complete
in G203. Packaging this application does not reserve it for a
later movement period. The remaining numerical mixed/noncritical
remainder bound, mixed theta/Selmer realization and H04.4 stay open.
No new question row or automatic gate closure is asserted.
